% Generated draft: requires mathematical and source review.
\documentclass[11pt]{article}
\usepackage[margin=25mm]{geometry}
\usepackage{fontspec}
\setmainfont{Latin Modern Roman}
\usepackage{amsmath,amssymb,mathtools,unicode-math}
\setmathfont{Latin Modern Math}
\usepackage{enumitem,needspace,xurl,hyperref}
\setlength{\parindent}{0pt}
\setlength{\parskip}{5pt}
\setlength{\emergencystretch}{3em}
\newcommand{\N}{\mathbb{N}}
\newcommand{\Z}{\mathbb{Z}}
\newcommand{\Q}{\mathbb{Q}}
\newcommand{\R}{\mathbb{R}}
\newcommand{\C}{\mathbb{C}}
\newcommand{\F}{\mathbb{F}}
\newcommand{\A}{\mathbb{A}}
\newcommand{\PP}{\mathbb{P}}
\newcommand{\E}{\mathbb{E}}
\newcommand{\eps}{\varepsilon}
\newcommand{\dd}{\,\mathrm{d}}
\begin{document}
\section*{UnsolvedMath ID: 10000013}
\textbf{Source problem code:} catalogue entry
\subsection{Definitions and mathematical statement}
Start with \(\mathcal Q_0=\{[0,1]^2\}\). At each step choose one square uniformly from the current collection, remove it and replace it by its four congruent quarter-squares. After \(n\) subdivisions there are \(1+3n\) current squares. This defines all stages on the same recursive probability space.

For a precise literal connected-union reading, let
\[
 D_n=\min\{|\mathcal A|:\mathcal A\subseteq\mathcal Q_n,
 \ \bigcup_{Q\in\mathcal A}Q\text{ is connected and contains }(0,0),(1,1)\}.
\]
The current squares here are closed, so corner contacts count as connectivity. If the source instead requires adjacency along a positive-length edge or another crossing convention, it defines a different quantity and must be specified.

The first question asks for a deterministic positive sequence \((a_n)\) such that \(D_n/a_n\) converges to a nondegenerate random variable. In the usual distributional interpretation the target is
\[
 \exists a_n>0\text{ deterministic}\ \exists D\text{ finite and not almost surely constant}:\quad
 D_n/a_n\ \Longrightarrow\ D.
\]
The excerpt does not specify the convergence mode; almost-sure or in-probability convergence would be stronger requests and are not asserted.

The second question asks whether a shortest crossing stabilizes as subdivision continues. A shortest square collection need not be unique, and the collections themselves change when their squares are subdivided. A precise stabilization notion must specify tie selection and whether it concerns a geometric path, coarse ancestral choices, or another projection of the changing crossings. No eventual equality of current-square lists or canonical path is silently imposed.

The supplied source-side and bounded modern audits report no complete proof or counterexample. The summary calls the notes ``2024'', while the supplied preserved copy has a 2020 archive timestamp; no publication date is inferred from that discrepancy. The open label remains a qualified source assessment rather than independently verified current status.
\subsection{Short English statement}
Randomly subdivide one current square into four at each step. Does the minimum square-count crossing between opposite corners have a deterministic nondegenerate scaling limit, and do shortest crossings stabilize in an appropriate sense? Connectivity, convergence and stabilization conventions need clarification.
\subsection{Sources}
\begin{itemize}
\item Local registry: \url{data/unsolvedmath/problems.json}, permanent UnsolvedMath ID \(10000013\), source code catalogue entry; supplied statement, background, and research summary, updated 2026-08-06.
\item Itai Benjamini, \emph{Coarse Geometry and Randomness}, Saint-Flour notes, Open problem 1.73, PDF page 16: \url{https://www.wisdom.weizmann.ac.il/~itai/stflouraug24.pdf#page=16}. Preserved copy supplied at \url{https://arquivo.pt/noFrame/replay/20201231041548id_/http://www.wisdom.weizmann.ac.il/~itai/stflouraug24.pdf}. Supplied text extraction checked against rendered PDF, accessed 2026-07-29.
\item The provenance reports cross-document deduplication, source-side review and a bounded modern audit with no complete resolution located; no separate resolution publication is supplied.
\end{itemize}
\end{document}
